\documentclass[12pt, twoside]{article}
\input{../preamble}

\fancyhead[LO]{La Scuola d'Italia / Dr.\ Huson / IB Math \\* 18 April 2026}
\fancyhead[RO]{\fbox{\begin{minipage}{6cm}\footnotesize First \& Last name:\\[0.5cm]\end{minipage}}}
\fancyfoot[C]{\thepage}

\begin{document}
\subsubsection*{5.14 Classwork: Logarithms and Exponentials (no calculator)}

\begin{enumerate}[itemsep=0.15cm]

%--- Problem 1: Product rule with exponential [5 marks] ---
%--- Source: AA SL 2022 May P1 TZ1 Q05 ---
\item[1.] [Maximum mark: 5]

Consider the curve with equation $y = (2x - 1)e^{kx}$,
where $x \in \mathbb{R}$ and $k \in \mathbb{R}$.

The tangent to the curve at the point where $x = 1$
is parallel to the line $y = 5xe^k$.

Find the value of $k$.

\begin{tikzpicture}
  \draw (-0.6,0) rectangle (11,13);
  \foreach \y in {1,2,...,12} \draw[dotted] (0,\y)--(10.5,\y);
\end{tikzpicture}

\newpage

%--- Problem 2: Logarithmic function [9 marks] ---
%--- Source: AA SL 2021 May P1 TZ2 Q05 ---
\item[2.] [Maximum mark: 9]

Consider the function $f$ defined by
$f(x) = \ln(x^2 - 16)$ for $x > 4$.

The following diagram shows part of the graph of $f$
which crosses the $x$-axis at point $A$, with coordinates
$(a, 0)$. The line $L$ is the tangent to the graph of $f$
at the point $B$.

\begin{center}
\begin{tikzpicture}[xscale=0.5, yscale=0.75] % Axes
  \draw[->] (-0.5,0) -- (20,0) node[right] {$x$};
  \draw[->] (0,-2) -- (0,7.5) node[above] {$y$};
  % Asymptote (separate from y-axis)
  \draw[dashed] (3.9,-2) -- (3.9,5);
  \node[above right] at (4,4.8) {$x = 4$};
  % Curve f(x) = ln(x^2 - 16) for x > 4
  \draw[thick, domain=4.02:18.5, samples=150]
    plot (\x, {ln(\x*\x - 16)});
  \node[above] at (15.5, 4) {$f$};
  % Point A (x = sqrt(17) ≈ 4.123)
  \node[circle, fill, inner sep=1.5pt] at (4.123, 0) {};
  \node[below right] at (4.123, -0.15) {$A(a, 0)$};
  % Point B (at x = 8, y = ln 48 ≈ 3.871)
  \node[circle, fill, inner sep=1.5pt] at (8, 3.871) {};
  \node[above left] at (7.9, 3.9) {$B$};
  % Tangent line L at B: y = (x-8)/3 + ln(48)
  \draw[thick, domain=5:17.5, samples=2]
    plot (\x, {(\x - 8)/3 + ln(48)});
  \node[right] at (15.5, 7) {$L$};
\end{tikzpicture}
\end{center}

\begin{enumerate}[itemsep=0.15cm]
  \item Find the exact value of $a$. \hfill [3]
  \item Given that the gradient of $L$ is $\dfrac{1}{3}$,
  find the $x$-coordinate of $B$. \hfill [6]
\end{enumerate}

\begin{tikzpicture}
  \draw (-0.6,0) rectangle (11,9.5);
  \foreach \y in {1,2,...,9} \draw[dotted] (0,\y)--(10.5,\y);
\end{tikzpicture}

\newpage

%--- Problem 3: Exponential function and inverse [14 marks] ---
%--- Source: Math SL 2019 Nov P1 Q10 ---
\item[3.] [Maximum mark: 14]

Let $g(x) = p^x + q$, for $x, \, p, \, q \in \mathbb{R}$, $p > 1$.
The point $A(0, a)$ lies on the graph of $g$.

Let $f(x) = g^{-1}(x)$. The point $B$ lies on the graph
of $f$ and is the reflection of point $A$ in the line $y = x$.

\begin{enumerate}[itemsep=0.15cm]
  \item Write down the coordinates of $B$. \hfill [2]
\end{enumerate}

The line $L_1$ is tangent to the graph of $f$ at $B$.

\begin{enumerate}[itemsep=0.15cm]
  \item[(b)] Given that $f'(a) = \dfrac{1}{\ln p}$, find the equation
  of $L_1$ in terms of $x$, $p$ and $q$. \hfill [5]
\end{enumerate}

The line $L_2$ is tangent to the graph of $g$ at $A$ and has
equation $y = (\ln p)\, x + q + 1$.

The line $L_2$ passes through the point $(-2, -2)$.

The gradient of the normal to $g$ at $A$ is
$\dfrac{1}{\ln \frac{1}{3}}$.

\begin{enumerate}[itemsep=0.15cm]
  \item[(c)] Find the equation of $L_1$ in terms of $x$. \hfill [7]
\end{enumerate}

\emph{Continue on lined paper.}

\end{enumerate}
\end{document}
